% Script section 6.3: ACK bug and slow convergence on a newer kernel.
\sectiontag{4}\source{}
\begin{frame}[t]{UC1: Linux Kernel TCP ACK Bug}
\definecolor{PacketBlue}{HTML}{0047CC}
\definecolor{OriginalPurple}{HTML}{7820A0}
\definecolor{AckGreen}{HTML}{007A35}
\definecolor{RetransmitRed}{HTML}{D02020}
\centering
\begin{tikzpicture}[x=1pt,y=1pt,
 every node/.style={font=\fontsize{20}{24}\selectfont\bfseries\boldmath,align=center},
 packet/.style={flow,draw=PacketBlue,line width=2.6pt},
 reply/.style={ack,draw=AckGreen,dashed,line width=2.6pt}]
\path[use as bounding box] (0,0) rectangle (890,285);
\node[anchor=east] at (108,210) {Sender};
\node[anchor=east] at (108,60) {Receiver};
\draw[-{Stealth[length=7pt]},SemGray,line width=1.2pt] (115,210)--(880,210);
\draw[-{Stealth[length=7pt]},SemGray,line width=1.2pt] (115,60)--(880,60);
\node[anchor=east,text=SemGray] at (880,135) {Time $\rightarrow$};
\draw[SemGray,densely dotted,line width=1.2pt] (220,230)--(220,42);
\node[text=SemGray] at (220,20) {RTT drop};
\node[anchor=south,text=OriginalPurple] at (145,245) {Original $P$};
\draw[OriginalPurple,line width=1.4pt] (145,244)--(140,214);
\draw[packet,draw=OriginalPurple,dashed] (140,210)--(550,60);
\node[text=OriginalPurple,fill=white,inner sep=3pt] at (195,163) {Delayed};
\node[anchor=south,text=PacketBlue] at (300,245) {Later packets};
\foreach \start/\finish in {245/295,280/330,315/365} {
  \draw[packet] (\start,210)--(\finish,60);
}
\uncover<2->{
\foreach \start/\finish in {295/345,330/380,365/415} {
  \draw[reply] (\start,60)--(\finish,210);
}
\node[text=AckGreen,anchor=north] at (330,44) {3 dup ACKs};
\node[anchor=south,text=RetransmitRed] at (460,245) {Fast retransmit $P$};
\draw[RetransmitRed,line width=1.4pt] (460,244)--(425,214);
\draw[packet,draw=RetransmitRed] (425,210)--(580,60);
}
\uncover<3->{
\fill[OriginalPurple] (550,60) circle (3pt);
\fill[RetransmitRed] (580,60) circle (3pt);
\node[anchor=north,text=OriginalPurple] at (545,44) {$P$};
\node[anchor=north,text=RetransmitRed] at (590,44) {Copy};
\draw[reply] (550,60)--(600,210);
\node[text=AckGreen,fill=white,inner sep=2pt,align=left] at (620,143) {ACK\\advances};
}
\uncover<4->{
\node[anchor=south,text=PacketBlue] at (705,245) {Further data};
\draw[packet] (640,210)--(690,60);
\draw[packet] (710,210)--(760,60);
\draw[reply] (690,60)--(740,210);
\draw[reply] (760,60)--(810,210);
\node[text=AckGreen,anchor=north] at (725,44) {ACK stays fixed};
\draw[RetransmitRed,line width=2.6pt] (850,202)--(850,218);
\node[anchor=south,text=RetransmitRed] at (850,232) {RTO};
}
\end{tikzpicture}\par
\vspace{8pt}
\raggedright
\begin{overlayarea}{\textwidth}{100pt}
\compactfont
\only<1>{High RTT for $\sim$1 s, then a sharp drop: later packets overtake $P$.\par}
\only<2>{Three duplicate ACKs trigger fast retransmit while the original $P$ is still in flight.\par}
\only<3>{The original $P$ arrives first and advances the ACK number. Its retransmitted copy then triggers the ACK-update bug.\par}
\only<4>{Further packets arrive, but the ACK number stays fixed; the sender eventually enters RTO.\par
\vspace{7pt}
{\labelfont Why BBR? Smooth pacing makes this arrival order easier to trigger.\par
\color{SemGray}Authors verified using eBPF and Wireshark.\par}}
\end{overlayarea}
\end{frame}

\sectiontag{4}\source{Fig. 6 (throughput panel), AdvNet.}
\begin{frame}[t]{UC1: BBR v3 Slow Convergence on a Newer Kernel}
\compactfont
On kernel 6.13.7+, the earlier ACK bug was no longer observed.\par
\vspace{18pt}
\includegraphics[width=\linewidth]{fig06_throughput.png}
\begin{center}\labelfont Time (seconds)\end{center}
\begin{itemize}
\setlength{\itemsep}{12pt}
\item CUBIC ramps up quickly; BBR v3 stays low over the 5-second run.
\pause
\item In a longer run, BBR v3 eventually ramps up: slow convergence, rather than a permanent stall.
\end{itemize}
\end{frame}

% Blueprint 18: fairness outcomes and the observed sending behavior.
\sectiontag{4}\source{}
\begin{frame}[t]{UC1: Fairness Between Competing Flows}
\compactfont
Compare the same flow's throughput alone ($R$) and with a competitor ($T$).\par
\vspace{8pt}
\begin{itemize}
\compactfont\setlength{\itemsep}{6pt}
\item CUBIC can take nearly all bandwidth and starve BBR.
\item In other environments, BBR can starve CUBIC.
\end{itemize}
\vspace{4pt}
\pause
\lead{Observed sending behavior}
\begin{tikzpicture}[x=1pt,y=1pt,
 every node/.style={font=\fontsize{20}{24}\selectfont,align=center},
 data/.style={flow,draw=blue!75!black,line width=2.2pt},
 feedback/.style={ack,draw=green!45!black,dashed,line width=2.2pt}]
\path[use as bounding box] (0,-8) rectangle (890,205);
\node[draw=SemBlue,line width=1.2pt,inner sep=9pt,text width=240pt] (normal) at (250,170) {Normal sending\\40--80 s};
\node[draw=SemBlue,line width=1.2pt,inner sep=9pt,text width=420pt] (limited) at (660,170) {One packet, then wait for its ACK\\Remaining 25--60 s};
\draw[flow,line width=2pt] (normal)--(limited);
\node[anchor=east] at (100,100) {Sender};
\node[anchor=east] at (100,10) {Receiver};
\draw[-{Stealth[length=7pt]},SemGray,line width=1.2pt] (110,100)--(880,100);
\draw[-{Stealth[length=7pt]},SemGray,line width=1.2pt] (110,10)--(880,10);
% Multiple packets are in flight during normal sending.
\foreach \send/\receive/\acktime in {145/200/250,180/235/285,215/270/320,250/305/355} {
  \draw[data] (\send,100)--(\receive,10);
  \draw[feedback] (\receive,10)--(\acktime,100);
}
\draw[SemGray,densely dotted,line width=1.2pt] (395,-3)--(395,120);
% After the transition, each new packet follows the previous packet's ACK.
\foreach \send/\receive/\acktime in {430/490/550,560/620/680,690/750/810} {
  \draw[data] (\send,100)--(\receive,10);
  \draw[feedback] (\receive,10)--(\acktime,100);
}
\node[text=blue!75!black,fill=white,inner sep=2pt] at (454,59) {Packet};
\node[text=green!45!black,fill=white,inner sep=2pt] at (525,49) {ACK};
\node[anchor=east,text=SemGray] at (880,55) {Time $\rightarrow$};
\end{tikzpicture}\par
\vspace{8pt}
{\labelfont\color{SemGray}The BBR team confirmed that the reported environment was problematic.\par}
\end{frame}

% Blueprint 19: distinguish fairness measurements from single-flow tests.
\sectiontag{4}\source{}
\begin{frame}[t]{UC2: Regression Testing Across BBR Versions}
\compactfont
\lead{Fairness regression (same trace)}
Reference ($R$): CUBIC throughput when competing with BBR v1.\par
Target ($T$): CUBIC throughput when competing with BBR v3.\par
\vspace{8pt}
{\centering
\begin{tikzpicture}[x=1pt,y=1pt,
 every node/.style={font=\fontsize{20}{24}\selectfont,align=center}]
\path[use as bounding box] (0,-28) rectangle (870,75);
\filldraw[fill=PaleBlue,draw=SemBlue] (0,0) rectangle (240,34);
\filldraw[fill=PaleTeal,draw=SemTeal] (240,0) rectangle (860,34);
\filldraw[fill=white,draw=SemTeal,line width=2pt] (757,0) rectangle (860,34);
\node at (120,17) {Dynamic trace};
\node at (485,17) {Fixed network conditions};
\node[font=\fontsize{18.5}{23}\selectfont] at (809,17) {Measure};
\draw[decorate,decoration={brace,amplitude=6pt},line width=1.2pt] (240,44)--(860,44);
\node at (550,65) {60 seconds};
\node[anchor=north] at (795,-3) {Final 10 s};
\end{tikzpicture}\par}
\vspace{8pt}
\pause
\begin{itemize}
\compactfont\setlength{\itemsep}{5pt}
\item With BBR v1: CUBIC gets $12\times$ BBR v1's throughput.
\item With BBR v3: BBR v3 gets $11.6\times$ CUBIC's throughput.
\end{itemize}
\vspace{10pt}
\lead{Separate single-flow test}
Without a competing CUBIC flow, AdvNet also finds traces where BBR v1 achieves more than $5\times$ BBR v3's throughput.\par
\end{frame}
